\subsection{window-\texorpdfstring{$6$}{6} 的 nongeometric replay core、boundary holonomy shell 与 repair-silent 刚性扇区}\label{subsec:conclusion-window6-replaycore-holonomy-boundary-blindspot}

沿用小节 \ref{subsec:conclusion-foldbin-stable-k0-faithful-boundary-parity-groupoid}、
\ref{subsec:conclusion-window6-mod2-k0-sync-replay-boundary-face}、
\ref{subsec:conclusion-window6-subcover-coinvariant-thermal-resonance-lock}、
\ref{subsec:conclusion-window6-boundary-spinor-shadow-gatejump-rigidity}、
定理 \ref{thm:conclusion-window6-local-index-binary-log-symmetry-surcharge}、
定理 \ref{thm:conclusion-window6-boundary-parity-zero-one-three-law}、
定理 \ref{thm:conclusion-window6-boundary-parity-minimal-torus-rank}、
定理 \ref{thm:conclusion-window6-minimal-shell-rigid-subcover-root-slice}、
定理 \ref{thm:window6-minimal-markov-obstruction-pure-f8-tail}、
定理 \ref{thm:window6-f8-repair-bit-no-boundary-covariant-extension}
与推论 \ref{cor:window6-f8-repair-vs-f9-boundary-parity-separated} 的记号。前述链条已经把 window-\(6\) 的裸可逆位数、boundary parity 的 \((0/1/3)\) 通道律、最小二重点层的 rigid \(X_4\)-子覆盖、以及 \(F_8/F_9\) 两条隐藏通道的分离固定下来。把这些接口继续并读之后，还可把同一 hidden sector 再压成一条更细的双层结构：一方面，最小精确回放的二位核心必然打破 canonical boundary parity，因此 \(b_{\mathrm{inv}}=2\) 真正度量的是 purely nongeometric replay core；另一方面，boundary parity 自身则作为一个三维 holonomy shell 独立存在，并在强局域几何与有理连续两类通道之外留下一个不可约的二维盲区。其间的 \(F_8/F_9\) 分裂又说明：repair 电荷与 boundary parity 虽同处最小壳附近，却分别占据不同的 hidden displacement 通道，不能再被归并为单一局域自由度。

\begin{theorem}[最小精确回放的对称破缺]\label{thm:conclusion-window6-minimal-exact-replay-symmetry-breaking}
记
\[
\pi:=\Fold_6^{\mathrm{bin}}:\Omega_6\to X_6.
\]
若 completion
\[
c:\Omega_6\to \{1,2,3,4\}
\]
满足
\[
(\pi,c):\Omega_6\longrightarrow X_6\times \{1,2,3,4\}
\]
为单射，则 \(c\) 不可能对 canonical boundary parity 中心
\[
Z_{\partial}\cong (\ZZ/2\ZZ)^3
\]
等变。等价地，任一实现 window-\(6\) 最小精确回放的二位解码器，都必然破坏 canonical boundary parity。
\end{theorem}
\begin{proof}
由定理 \ref{thm:conclusion-window6-local-index-binary-log-symmetry-surcharge}，
\[
b_{\mathrm{inv}}=2,
\]
即 \(B=2\) 已足以实现精确反演。若存在四值 completion \(c\) 使 \((\pi,c)\) 单射且 \(c\) 同时对 \(Z_{\partial}\) 等变，则这正给出一个 boundary-equivariant 的四值伴随场。然同一定理与推论 \ref{cor:xi-time-part9x-no-boundary-equivariant-four-valued-companion-field} 已经排除：任何保持 canonical boundary parity 的 completion 都必须满足
\[
B_{\partial}\ge 3,
\]
因而四值标签不可能与 boundary 等变性并存。故所有 \(B=2\) 的最小精确回放方案都必然打破 canonical boundary parity。证毕。
\end{proof}

\begin{theorem}[boundary 二挠的双重最小实现]\label{thm:conclusion-window6-boundary-twofold-minimal-faithful-realization}
window-\(6\) 的 canonical boundary parity 具有两种彼此独立而同秩的最小忠实实现：
\[
\operatorname{rank}_{\FF_2} Z_{\partial}=3,
\qquad
r_{\mathrm{tor}}(Z_{\partial})=3.
\]
并且 \(Z_{\partial}\) 在交换化电荷格
\[
\bigl(G_6^{\mathrm{bin}}\bigr)^{\mathrm{ab}}\cong (\ZZ/2\ZZ)^{21}
\]
中给出一条规范的三维坐标子格。
\end{theorem}
\begin{proof}
由推论 \ref{cor:fold-bin-boundary-center-z2}，三条 boundary 二点纤维各自产生一个独立的中心 sheet-flip，故
\[
Z_{\partial}\cong (\ZZ/2\ZZ)^3,
\qquad
\operatorname{rank}_{\FF_2} Z_{\partial}=3.
\]
再由推论 \ref{cor:conclusion-window6-anomaly-root-torus-mod2-splitting}，
\[
\bigl(G_6^{\mathrm{bin}}\bigr)^{\mathrm{ab}}
\cong
\FF_2^{X_6^{\mathrm{cyc}}}\oplus \FF_2^{X_6^{\mathrm{bdry}}},
\]
其中
\[
\FF_2^{X_6^{\mathrm{bdry}}}\cong Z_{\partial}.
\]
故 \(Z_{\partial}\) 在交换化电荷格中恰给出一条规范三维坐标子格。

另一方面，由定理 \ref{thm:conclusion-window6-boundary-torus-2torsion}，
\[
Z_{\partial}\cong T_6[2],
\qquad
T_6\cong (S^1)^3.
\]
再由定理 \ref{thm:conclusion-window6-boundary-parity-minimal-torus-rank}，faithful 实现全部 boundary parity 所需的最小 torus holonomy 秩正是 \(3\)。故 boundary parity 在离散交换化语义与紧环面 holonomy 语义下同时达到同一最小忠实秩。证毕。
\end{proof}

\begin{theorem}[hidden displacement 谱的 \texorpdfstring{$(21/34)$}{(21/34)} 分裂]\label{thm:conclusion-window6-hidden-displacement-21-34-splitting}
在 window-\(6\) 的审计位移谱
\[
\Delta_6=\{0,21,34,55\}
\]
中，最小一阶 Markov 修复通道严格承载于
\[
21=F_8,
\]
而一切保持稳定类型标签且由强局域 geometric uplift 实现的非平凡隐藏位移，严格承载于
\[
34=F_9.
\]
因此，Markov 修复通道与几何可实现的 boundary parity 通道在 hidden displacement 谱上已发生严格分裂；二者既不可互代，也不能由同一局域几何机制统一生成。
\end{theorem}
\begin{proof}
由定理 \ref{thm:window6-minimal-markov-obstruction-pure-f8-tail}，window-\(6\) 的最小一阶 Markov 修复完全落在被截断的 \(F_8\) 尾位通道，即
\[
\delta=21=F_8,
\]
并且这一方向不属于任何几何二点方向谱。

另一方面，由推论 \ref{cor:terminal-window6-local-uplift-admissibility} 以及结论 \ref{con:xi-terminal-window6-local-uplift-so10-unique}，在保持稳定类型标签且由 \(\Aut(Q_6)\) 实现的强局域 geometric uplift 口径下，除平凡位移外仅允许
\[
\delta=\pm 34,
\]
亦即唯一局域可实现的非零隐藏位移严格位于 \(F_9\) 通道。

再由定理 \ref{thm:window6-f8-repair-bit-no-boundary-covariant-extension} 与推论 \ref{cor:window6-f8-repair-vs-f9-boundary-parity-separated}，
\[
b_6^{\mathrm{bin}}|_B\equiv 0,
\]
故 \(F_8\) 修复位在 boundary 模块上既不承担反对称 sheet-flip，也不可能与 canonical \(F_9\)-boundary parity 统一。于是 \((21)\)-修复通道与 \((34)\)-几何通道必然分裂。证毕。
\end{proof}

\begin{theorem}[最小回放位数、几何余核维数与 cubical 秩的三重同一]\label{thm:conclusion-window6-replay-core-geometry-defect-cubical-rank-identity}
定义 window-\(6\) 的局域 cubical 秩
\[
r_{\mathrm{cube}}
:=
\min\Bigl\{
r:\ \forall w\in X_6,\ F_w
\text{ 可仿射同构为 }\{0,1\}^r\text{ 的某个序理想}
\Bigr\}.
\]
则有严格等式
\[
b_{\mathrm{inv}}
=
\dim_{\FF_2}(P_6/P_6^{\mathrm{geo}})
=
r_{\mathrm{cube}}
=2,
\]
其中
\[
P_6:=Z_{\partial}\cong (\ZZ/2\ZZ)^3,
\qquad
P_6^{\mathrm{geo}}=\langle(1,1,1)\rangle\cong \ZZ/2\ZZ.
\]
因此，window-\(6\) 的 nongeometric replay core 恰由两位构成：它一方面等于最小精确回放预算，另一方面又等于 geometric quotient 的残余维数，并与 hidden fiber 的最小 cubical 维数完全一致。
\end{theorem}
\begin{proof}
第一式由定理 \ref{thm:conclusion-window6-local-index-binary-log-symmetry-surcharge} 给出：
\[
b_{\mathrm{inv}}=2.
\]
第二式由定理 \ref{thm:conclusion-window6-boundary-parity-zero-one-three-law} 中的几何部分得到
\[
P_6^{\mathrm{geo}}=\langle(1,1,1)\rangle,
\qquad
P_6/P_6^{\mathrm{geo}}\cong (\ZZ/2\ZZ)^2,
\]
故
\[
\dim_{\FF_2}(P_6/P_6^{\mathrm{geo}})=2.
\]

对第三式，定理 \ref{thm:xi-time-part9xab-window6-hidden-fiber-affine-boolean-normal-form} 已证明：对每个 \(w\in X_6\)，hidden fiber 都具有标准形
\[
F_w=V_6(w)+\beta(I_{d(w)}),
\qquad
\beta(a,b)=21a+34b,
\]
其中 \(I_2,I_3,I_4\subset \{0,1\}^2\) 都是 \(\{0,1\}^2\) 的序理想。故
\[
r_{\mathrm{cube}}\le 2.
\]
另一方面，window-\(6\) 存在大小为 \(4\) 的纤维，而任意 \(\{0,1\}^1\) 的序理想至多只含 \(2\) 个点，故
\[
r_{\mathrm{cube}}\ge 2.
\]
于是
\[
r_{\mathrm{cube}}=2.
\]
三式合并即得结论。证毕。
\end{proof}

\begin{proposition}[boundary parity 的 rigid \texorpdfstring{$X_4$}{X4} 局域化]\label{prop:conclusion-window6-boundary-parity-rigid-x4-localization}
记
\[
S_{6,2}:=\{w\in X_6:\ d_6^{\mathrm{bin}}(w)=2\},
\qquad
\iota_4:X_4\hookrightarrow X_6,
\quad
\iota_4(v)=(v,0,1).
\]
则
\[
S_{6,2}=\iota_4(X_4),
\qquad
X_6^{\mathrm{bdry}}\subset S_{6,2}.
\]
从而 canonical boundary parity 并不分布在一般可见态之上，而是完全局域在一个 rigid、余维二且后缀锁定的 lower-resolution stable subcover 上。
\end{proposition}
\begin{proof}
定理 \ref{thm:conclusion-window6-minimal-shell-rigid-subcover-root-slice} 已证明
\[
S_{6,2}=\iota_4(X_4).
\]
而同一定理又给出
\[
X_6^{\mathrm{bdry}}
=
\{\texttt{100001},\texttt{100101},\texttt{101001}\}
\subset S_{6,2}.
\]
因此 boundary parity 的全部载体都嵌入 \(\iota_4(X_4)\) 这一刚性子覆盖中。证毕。
\end{proof}

\begin{theorem}[repair-silent boundary module]\label{thm:conclusion-window6-repair-silent-boundary-module}
设
\[
B:=\bigcup_{w\in X_6^{\mathrm{bdry}}}F_w
\]
为 boundary 三条二点纤维之并。则在 boundary 模块 \(B\) 上，canonical boundary parity 为非平凡反对称二挠，而最小一阶 Markov 修复位满足
\[
b_6^{\mathrm{bin}}|_B\equiv 0.
\]
因此，window-\(6\) 的 boundary 模块是一个 parity-active 而 repair-silent 的刚性 hidden sector。
\end{theorem}
\begin{proof}
由推论 \ref{cor:fold-bin-boundary-center-z2}，三条 boundary 二点纤维的纤维内换位生成规范中心子群
\[
Z_{\partial}\cong (\ZZ/2\ZZ)^3,
\]
故 canonical boundary parity 在 \(B\) 上非平凡，并以反对称方式作用于每条 boundary 二点纤维。

另一方面，定理 \ref{thm:window6-f8-repair-bit-no-boundary-covariant-extension} 已证明
\[
b_6^{\mathrm{bin}}|_B\equiv 0.
\]
于是 boundary 模块上的拓扑二挠与一阶 Markov 修复电荷严格解耦：前者 active，后者 silent。证毕。
\end{proof}

\begin{corollary}[有理连续与强局域几何的联合盲区]\label{cor:conclusion-window6-rational-geometric-joint-blind-spot}
在
\[
P_6\cong (\ZZ/2\ZZ)^3
\]
内部，任一有理连续证书只能探测平凡子群，而任一强局域 geometric uplift 只能保留对角子群
\[
P_6^{\mathrm{geo}}\cong \ZZ/2\ZZ.
\]
故商群
\[
P_6/P_6^{\mathrm{geo}}\cong (\ZZ/2\ZZ)^2
\]
同时逃逸于有理连续包络与强局域几何两类机制之外；这给出一个不可约的余维二 parity 盲区。
\end{corollary}
\begin{proof}
由定理 \ref{thm:conclusion-window6-boundary-parity-zero-one-three-law}，
\[
r_{\mathrm{rat}}=0,
\qquad
r_{\mathrm{geo}}=1,
\qquad
P_6^{\mathrm{geo}}=\langle(1,1,1)\rangle.
\]
因此，有理连续包络对全部 boundary parity 完全失显，而强局域 geometric uplift 只保留唯一对角方向。遂得
\[
P_6/P_6^{\mathrm{geo}}\cong (\ZZ/2\ZZ)^2
\]
中的两位残余既非有理连续可见，也非强局域几何可实现。证毕。
\end{proof}

\noindent
这些结论把前一轮的 boundary-spinor/gate-jump 链继续压成一个更细的双层结构：\(b_{\mathrm{inv}}=2\) 不再只是一个局部预算，而是与 geometric quotient 的二维余核、以及 hidden fiber 的二维 cubical 标准形完全同一，从而刻画出一个纯粹的 nongeometric replay core；与之相对，\(Z_{\partial}\cong (\ZZ/2\ZZ)^3\cong T_6[2]\) 则给出独立的 boundary holonomy shell。二者在 \(F_8/F_9\) hidden displacement 谱上进一步分裂成 repair 与 parity 两条通道：前者在 boundary 模块上完全 silent，后者则完全局域于 rigid 的 \(X_4\)-子覆盖。最终留下的
\[
P_6/P_6^{\mathrm{geo}}\cong (\ZZ/2\ZZ)^2
\]
正是同时逃逸于有理连续包络与强局域几何的最小不可约盲区。
